\chapter{Literature Review and Theoretical Background}

\section{Introduction}
This fast development of deep learning and computer vision has changed the way maritime surveillance research is conducted, and a series of studies on intelligent, automated solutions to traditional radar and manual surveillance systems have emerged. The chapter reviews the literature on the three main technical aspects of the proposed system, namely, ship detection, multi-object tracking, and trajectory prediction, from the latest publications in the most prestigious journals such as Expert Systems with Applications, Engineering Applications of Artificial Intelligence, Journal of Marine Science and Engineering and Scientific Reports. The characteristics of each reviewed work are analyzed for their major contributions and methodological decisions as well as their limitations, and the works are specifically analyzed in terms of how well they address the complexities of the real world in maritime surveillance. The gaps as a whole that have been identified from literature are the technical and operational justification for the end-to-end integrated system proposed in this project.

\subsection{Deep Learning-Based Ship Detection}
Ship detection from optical and \gls{sar} imagery has been one of the most actively studied problems in maritime remote sensing over the past decade. The following subsections review contributions spanning improved \gls{yolov8} architectures, lightweight variants, multi-scale fusion strategies, and cross-domain adaptations that collectively illuminate the state of the art against which the present system is positioned.
 
\subsection{\textbf{\texorpdfstring{\gls{yolov8}}{YOLOv8} Variants for \texorpdfstring{\gls{sar}}{SAR} and Optical Ship Detection}}
 
Ning et al.\cite{Ning2024} introduced a \gls{yolov8}-oriented bounding-box (\gls{obb}) rotating detection algorithm tailored for \gls{sar} ship targets. By replacing the conventional axis-aligned bounding box regression head with an oriented prediction head, the method addressed the inherent difficulty of detecting arbitrarily rotated vessels in \gls{sar} scenes where clutter and speckle noise significantly degrade detection quality. The authors reported that oriented boxes reduced false positives caused by background interference near harbour areas. This work is directly relevant to the present project in that it establishes the backbone utility of \gls{yolov8} for maritime targets and demonstrates that architectural modifications to the detection head yield measurable precision gains, a principle echoed in the confidence-threshold-driven detection module employed in the proposed system.

\setlength{\parskip}{1em}
 
Gao et al. \cite{Gao2024} proposed SM-LightNet, a lightweight ship detection framework built on an improved \gls{yolov8} backbone. The core motivation was deployment on resource-constrained edge platforms where full \gls{yolov8} inference rates are insufficient for real-time operation. The authors introduced depthwise separable convolutions and a streamlined feature pyramid to reduce the parameter count without sacrificing mean average precision (\gls{map}) on standard maritime benchmarks. The reported inference speeds demonstrated real-time capability on a mobile \gls{gpu}, establishing that computational efficiency and detection accuracy are not mutually exclusive objectives. The emphasis on edge deployment resonates with the \gls{cuda}-accelerated inference strategy adopted in the proposed pipeline, where model selection between \gls{yolov8}n and \gls{yolov8}m was governed by the same speed-accuracy trade-off.

\setlength{\parskip}{1em}

 
Gong et al.\cite{Gong2024} presented a real-time long-distance ship detection architecture based on \gls{yolov8}, published in IEEE Access. The key challenge addressed was the degradation of detection performance when vessels occupy a very small fraction of the image area due to camera range. The authors augmented the standard \gls{yolov8} feature pyramid with an additional super-resolution branch that reconstructed fine-grained spatial detail for distant targets. Evaluations on a proprietary long-range maritime dataset demonstrated that the enhanced architecture recovered detection recall for ships as small as $16 \times 16$ pixels. The relevance to the proposed work lies in the shared challenge of detecting vessels that may appear at varied scales depending on camera altitude and zoom, which the system addresses
by retaining \gls{yolov8}'s native multi-scale prediction heads unchanged while relying on the drone stabiliser to maintain a consistent apparent scale across frames.

\setlength{\parskip}{1em}
 
Chen and Liu \cite{Chen2024} focused specifically on infrared ship target detection, adapting \gls{yolov8} with a custom attention module to exploit the distinct thermal signatures of vessels in low-visibility infrared imagery. The study demonstrated that standard RGB-optimised backbone features were sub-optimal for infrared input, motivating the insertion of a channel-wise recalibration block that reweighted feature map activations according to their infrared discriminability. Results on a curated infrared maritime dataset showed substantially improved detection of small, dim targets in cold-sea backgrounds. Although the present project operates on optical video rather than infrared, the insight that domain-specific feature recalibration improves detection of spectrally atypical targets informs the rationale for fine-tuning the \gls{yolov8} backbone on maritime-specific imagery.
 
Yu and Shin \cite{Yu2023} investigated the adaptation of \gls{yolov8} for \gls{sar}-based ship detection, focusing on the fundamental mismatch between the natural image statistics on which \gls{yolov8} was originally trained and the coherent-imaging characteristics of \gls{sar} data. The authors introduced a preprocessing normalisation pipeline that standardised \gls{sar} amplitude values into a distribution compatible with \gls{yolov8}'s convolutional filters, yielding consistent detection performance across different \gls{sar} sensors and incidence angles. Their ablation study confirmed that input normalisation contributed more to final \gls{map} than architectural modifications alone, underscoring the importance of data preprocessing in cross-domain detection. This finding guided the optical preprocessing choices in the proposed system, specifically the stabilisation stage that normalises spatial reference frames prior to detection.
 
\subsection{\textbf{Lightweight and Efficient Ship Detection Architectures}}
 
Liu et al.\cite{Liu2025} addressed ship target detection in high-resolution remote sensing imagery by refining the \gls{yolov8} neck architecture with a bidirectional feature pyramid that improved the propagation of both fine-grained and semantic features across scales. The authors reported competitive detection rates on the \gls{dota} and HRSC2016 benchmarks while constraining the model's floating point operations per second (\gls{flops}) to within the budget of an onboard satellite processor. The paper's discussion of the challenge of intra-class scale variance,
where cargo ships and fishing vessels may differ by two orders of magnitude in pixel area, directly parallels the vessel scale variation encountered in the drone-mounted camera setup of the present project.
 
\setlength{\parskip}{1em}

Fu et al.\cite{Fu2025} extended the infrared ship detection thread by proposing a \gls{yolov8}-based network augmented with a deformable convolutional layer to handle irregular thermal silhouettes that arise when vessels are partially obscured by sea clutter or atmospheric haze. The deformable kernels allowed the receptive field to conform to the elongated, non-rectangular shapes typical of ship superstructures
in infrared imagery. Experiments on the \gls{flir} maritime dataset recorded a mean detection precision gain of over five percentage points compared to baseline \gls{yolov8}. The demonstrated benefit of deformable receptive fields for irregular targets reinforces the complementary role of the \gls{osnet} \gls{reid} module in the proposed tracker, which similarly adapts its feature extraction to non-rigid appearance deformations caused by viewing angle changes.

\setlength{\parskip}{1em}

Yuan et al. \cite{Yuan2024} proposed EN-\gls{yolo}, an enhanced ship detection network
that augmented the \gls{yolov8} backbone with an explicit edge-detection branch running
in parallel with the main feature pyramid. The additional branch produced edge-map
feature tensors that were concatenated with the standard convolutional outputs
before the detection head, effectively providing the network with structural
boundary cues that are especially informative for watercraft targets whose hull
outlines are well-defined against open-water backgrounds. Benchmark results on
the SeaShips dataset showed a 3.8\% \gls{map} improvement over the \gls{yolov8} baseline.
The principle that explicit structural priors complement learned semantic features
is consistent with the bounding-box regression approach adopted in the proposed
system, where \gls{yolov8}'s built-in anchor-free head provides precise hull
localisation from which centroid coordinates are derived.
 
Shi et al. \cite{Shi2025} proposed a remote sensing ship target recognition
algorithm using an improved \gls{yolov8} with a convolutional block attention module
(\gls{cbam}) integrated into each C2f block of the backbone. The attention mechanism
enforced suppression of irrelevant sea-clutter channels while amplifying
ship-relevant spatial regions, producing a cleaner feature map entering the
detection head. Evaluation on the \gls{sar}-Ship dataset yielded recall rates exceeding
92\%, with particular improvements on small targets (area $<$ 64~px\textsuperscript{2}).
The demonstrated efficacy of attention-guided feature suppression for small
maritime targets is relevant to the proposed project's handling of distant vessels
that may subtend only tens of pixels in the drone camera's field of view.
 
\subsubsection{\textbf{Attentive and Multi-Scale Detection for \gls{sar} Imagery}}
 
Iliger et al. \cite{Iliger2024} presented an attentive \gls{yolov8} architecture for
\gls{sar} ship detection in which a transformer-based self-attention block was inserted
between the backbone and the neck. The self-attention layer produced
globally-informed feature tokens that captured long-range spatial dependencies
between candidate ship locations and their surrounding context, enabling the
network to suppress false alarms in densely clustered port environments. The
study evaluated the approach on the \gls{ssdd} and \gls{sar}-Ship datasets, reporting
significant reductions in false positive density at busy anchorage zones. The
relevance to the present work is methodological: the global context modelling
principle, here implemented via transformers, is approximated in the proposed
pipeline through the 90-frame track-history buffer maintained per vessel, which
provides temporal context equivalent to the spatial context captured by the
attention mechanism.
 
Wang et al. \cite{Wang2025} studied ship detection in satellite imagery using
\gls{yolov8} improved with a spatial pyramid pooling fast (\gls{sppf}) module extended to
five pooling scales instead of the standard three. The enlarged receptive field
coverage was shown to improve detection of large cargo ships that spanned multiple
\gls{yolov8} feature pyramid levels, a category that standard architectures tend to
localise poorly because the ship's bow and stern may lie in separate anchor
regions. The reported \gls{map} gain of 4.2\% on the \gls{dota}-v1.5 ship category
demonstrated that multi-scale pooling is an effective remedy for this large-object
fragmentation problem. The proposed system addresses an analogous fragmentation
risk by using the full bounding box (rather than the centroid alone) as the
primary detection primitive passed to the tracker, thereby preserving large-vessel
spatial extent information.
 
Lin and Han \cite{Lin2024} introduced GCSAR, a gradient-aggregation and
channel-enhancement algorithm for small-object \gls{sar} ship detection. The key
innovation was a gradient-guided feature aggregation module that combined
multi-level gradient maps with convolutional feature tensors to emphasise
spatially fine-grained ship signatures. Channel enhancement was applied through
a lightweight squeeze-and-excitation block that redistributed feature map energy
toward the most discriminative channels. Experiments on the \gls{ssdd} dataset
demonstrated improvements specifically on the smallest size class of vessels,
a category historically underserved by general-purpose detectors. The problem
of small vessel detection is directly applicable to the proposed system's
deployment scenario, where vessels viewed from elevated drone altitude may
register as small-area bounding boxes.
 
Huang et al. \cite{Huang2024} explored enhanced multi-scale ship detection
through joint training on multi-source image data, combining optical, \gls{sar}, and
infrared imagery within a single training pipeline using domain-adaptive batch
normalisation layers. The cross-domain training strategy produced a feature
extractor that generalised across imaging modalities, achieving competitive
performance on each modality individually while requiring only a single model
at inference time. The demonstrated benefit of multi-source data augmentation
for learning robust features is consistent with the data augmentation strategies
employed during \gls{yolov8} fine-tuning in the proposed system, which included
brightness jitter, mosaic augmentation, and horizontal flipping to simulate
diverse maritime lighting and viewpoint conditions.
 
Chen et al. \cite{Chen2023} presented a real-time ship detection algorithm based
on an improved \gls{yolov8} network, targeting surveillance video from coastal camera
systems. The authors modified the \gls{yolov8} decoupled head to share more parameters
between the classification and regression branches, reducing inference latency by
approximately 12\% without measurable \gls{map} degradation. A dynamic non-maximum
suppression (\gls{nms}) algorithm was also introduced to reduce the post-processing
overhead for densely populated frames. The paper reported processing rates
exceeding 30 frames per second on mid-tier \gls{gpu} hardware, validating the suitability
of \gls{yolov8} for continuous-stream video surveillance. These latency benchmarks
directly informed the hardware specification choices for the proposed system and
confirmed that real-time throughput at the target frame rate was achievable with
the selected \gls{yolov8}m variant.
 
\subsubsection{\textbf{Remote Sensing and Satellite-Based Ship Detection}}
Yu et al. \cite{Yu2026} proposed a \gls{yolov8}-based system for ship detection in
general surveillance video, emphasising the challenge of handling partial
occlusion caused by bridge structures and harbour infrastructure that frequently
intersect with vessel bounding boxes. The authors introduced an occlusion-aware
loss function that penalised false negatives arising from partially visible ships
more heavily than those from fully unoccluded ships, producing a detector with
higher recall under occlusion at a marginal precision cost. The reported recall
improvement of 7.3\% on an occlusion-heavy test split is directly relevant to the
present system, which similarly encounters vessels that pass behind harbour
structures or are partially blocked by other ships in a congested scene.
 
Lee et al. \cite{Lee2025} developed a lightweight improved \gls{yolov8} for \gls{sar}-based
ship detection aboard \gls{leo} satellites within Integrated
Communication and Remote Sensing (\gls{icrs}) systems. Given the extreme constraints
on onboard processing power and uplink bandwidth, the authors pruned \gls{yolov8}
to approximately 30\% of its original parameter count using a structured magnitude
pruning strategy, then recovered accuracy through knowledge distillation from the
full-size teacher model. The resulting student model achieved inference within
the power budget of a CubeSat-class onboard computer while retaining 96\% of
the teacher's \gls{map}. Although the proposed surveillance pipeline does not operate
on satellite hardware, the knowledge distillation methodology represents a
general model compression technique applicable to future edge-deployment variants
of the system.
 
Cao et al. \cite{Cao2024} proposed an improved \gls{yolov8} algorithm for ship
detection with a focus on the challenge of inshore clutter, where ships berthed
at docks produce detection responses that overlap with harbour infrastructure.
The authors introduced a context-aware suppression module that examined the pixel
neighbourhood around each predicted bounding box and reduced confidence scores
for detections whose surrounding context matched known dock or pier texture
patterns. This approach reduced false-positive detections in port areas by 28\%
relative to baseline \gls{yolov8}. The challenge of distinguishing berthed from underway
vessels is directly analogous to the present system's need to distinguish
stationary vessels from moving ones in the tracking stage, which is resolved through
the stationary-check threshold applied before trajectory prediction is initiated.
 
Liu et al. \cite{Liu2025a} presented an optimised ship detection algorithm
for load test data mining using \gls{yolov8}. The work addressed the specific challenge
of detecting ships in noisy, low-contrast imagery captured under adverse weather
conditions during maritime load-testing operations. The authors applied a test-time
augmentation strategy that averaged predictions over multiple transformed versions
of each input frame, achieving a consistent \gls{map} improvement of approximately 3\%
across all weather categories. The principle of inference-time data augmentation
for robustness to image degradation is aligned with the proposed system's
optical-flow stabilisation stage, which similarly mitigates a specific form of
image quality degradation --- namely motion blur caused by platform jitter ---
before detection is performed.
 
Deshmukh and Joshi \cite{Deshmukh2025} examined the transition from semantic
segmentation-based ship detection methods to \gls{yolov8}, providing a systematic
comparison of the two paradigms on a shared maritime dataset. The study
demonstrated that while semantic segmentation produced more precise ship masks,
its inference latency was three to five times higher than \gls{yolov8} and its
performance degraded more sharply under partial occlusion, because segmentation
masks could not be reliably predicted for the hidden portions of a vessel.
\gls{yolov8}'s bounding-box representation proved more robust to occlusion precisely
because the box provides a complete localisation prior regardless of visibility.
This comparative finding validates the bounding-box detection strategy of the
proposed pipeline over mask-based alternatives for real-time surveillance
applications.
 
\subsubsection{\textbf{Specialised and Novel Detection Architectures}}
 
Liu and Yang \cite{Liu2025b} investigated the \gls{yolov8}n variant, the smallest
member of the \gls{yolov8} family, and evaluated its suitability for ship detection
tasks. The authors quantified the accuracy-efficiency frontier of \gls{yolov8}n
relative to \gls{yolov8}s, \gls{yolov8}m, and \gls{yolov8}l on the SeaShips and SMD datasets,
establishing empirical guidelines for model selection under different deployment
constraints. Their findings showed that \gls{yolov8}n was competitive with \gls{yolov8}s
in recall but lagged in precision for small targets, motivating the use of
\gls{yolov8}m as a default choice when moderate computational resources are available.
These comparative benchmarks directly informed the model selection rationale
in the proposed system, where \gls{yolov8}m was selected for the primary detection
pipeline and \gls{yolov8}n was retained as a configuration option for edge deployment.
 
Hong et al. \cite{Hong2024} introduced an adaptive spatial feature fusion
algorithm for \gls{sar} ship detection that dynamically weighted feature contributions
from different spatial scales based on the estimated ship size at each detection
location. A lightweight regression sub-network predicted the probable ship area
from local context, and this estimate was used to bias the feature pyramid fusion
coefficients accordingly. The approach produced consistent improvements across all
ship size categories on the \gls{hrsid} dataset, eliminating the traditional accuracy
gap between small and large target detection. The adaptive feature weighting
concept is analogous to the speed-scaled prediction window employed in the
proposed trajectory forecasting module, where the number of prediction steps
is dynamically adjusted based on the estimated vessel speed.
 
Zheng et al. \cite{Zheng2024} proposed a lightweight \gls{sar} ship target detection
algorithm by pruning and reshaping the \gls{yolov8} backbone with ghost modules, which
replaced standard convolutions with a combination of cheap linear operations and
a small number of intrinsic feature operations. The resulting Ghost\gls{yolov8} model
achieved a 40\% reduction in parameters and a $2.1\times$ inference speedup
compared to \gls{yolov8}n while maintaining competitive \gls{map} on the \gls{ssdd} benchmark.
The work also analysed the effect of different ghost ratios on the trade-off
between efficiency and accuracy, providing a principled methodology for tuning
lightweight architectures. These efficiency gains demonstrate the practical
feasibility of deploying powerful detection architectures on constrained hardware,
reinforcing the real-time viability of the proposed system.
 
Jiang et al. \cite{Jiang2024} proposed \gls{yolov8}-FDF, a small target detection
algorithm that introduced a feature distinction and fusion (FDF) module into the
\gls{yolov8} neck. The FDF module separately processed high-frequency (edge and texture)
and low-frequency (shape and region) components of feature maps, then fused them
with a learned attention gate that emphasised high-frequency cues for small
targets while retaining low-frequency context for large ones. Published in IEEE
Access, the study demonstrated significant \gls{map} improvements on custom complex
scene datasets and contributed a systematic analysis of how feature frequency
decomposition benefits detection of scale-diverse targets. The frequency-domain
perspective on feature richness complements the multi-scale detection heads
built into \gls{yolov8} and used in the proposed system.
 
Zhu et al. \cite{Lejun2024} introduced DGSP-\gls{yolo}, a high-precision \gls{sar} ship
detection model built on \gls{yolov8} with a dual global-sensing path. The global
sensing path employed dilated convolutions with multiple dilation rates to
capture multi-scale context without increasing the spatial downsampling ratio,
preserving fine-grained ship texture across the full feature hierarchy.
Published in IEEE Access, the model outperformed prior \gls{sar} detectors including
YOLOX and RT-DETR on the \gls{ssdd} and \gls{sar}-Ship benchmarks. The dilated convolution
strategy effectively expanded the effective receptive field without parameter
cost, a technique with direct application to the neck architecture of future
versions of the proposed system where broader context could improve detection
of convoy-formation vessels.
 
Dong et al. \cite{Dong2025} developed a component model for large ship detection
in \gls{sar} images, recognising that large vessels are structurally decomposable into
distinct sub-parts --- hull, superstructure, funnel --- whose relative arrangements
could be used as a complementary detection cue. The component model first located
individual parts with a lightweight detector, then used a geometric consistency
check to aggregate part detections into ship-level predictions. Published in the
IEEE Journal of Selected Topics in Applied Earth Observations and Remote Sensing,
the approach achieved state-of-the-art performance on the LS-\gls{ssdd}-v1.0 large-ship
dataset. While the proposed system does not implement a component model, the
principle of exploiting sub-part geometry for vessel-level reasoning may be
relevant for future extensions targeting vessel type classification.
 
Xu et al. \cite{Xu2025} presented RLE-\gls{yolo}, a lightweight and multi-scale \gls{sar}
ship detection network based on \gls{yolov8}, incorporating a re-parameterised linear
encoder (RLE) module that fused multi-branch convolutional responses into a single
branch at inference time. This structural re-parameterisation preserved the
representational diversity of the multi-branch training topology while reducing
inference cost to that of a single-branch network. Published in IEEE Access, the
model demonstrated strong performance on both the \gls{ssdd} and \gls{sar}-Ship datasets with
approximately 30\% fewer \gls{flops} than baseline \gls{yolov8}. The technique is broadly
applicable to the proposed system's \gls{yolov8} backbone and represents a candidate
optimisation for future deployment on lower-power platforms.
 
Qiu and Zhang \cite{Qiu2026} presented OMAP2-\gls{yolo}, a ship target detection
algorithm for \gls{sar} images addressing multi-scale ships in complex environments,
published in IEEE Transactions on Aerospace and Electronic Systems. The authors
introduced an orientation-aware multi-scale attention pyramid that simultaneously
estimated ship orientation and refined detection across five scale levels, directly
handling the problem of elongated large ships whose aspect ratio causes them to
straddle multiple anchor regions. Evaluation on the LS-\gls{ssdd} and \gls{hrsid} datasets
showed superior detection rates for both small fishing boats and large tankers
within a single unified model. The orientation estimation component is directly
relevant to the heading computation module in the proposed system, which derives
vessel heading from centroid displacement sequences.
 
Yang et al. \cite{Yang2024} developed \gls{yolo}-RDFEA for object detection in
\gls{rd} imagery using an improved \gls{yolov8} with feature enhancement
and attention mechanisms, published in IEEE Access. \gls{rd} imagery presents distinct
challenges compared to optical or \gls{sar} imagery because ship returns appear as
smeared range-Doppler signatures whose extents depend on vessel speed and radar
parameters. The authors addressed this with a Doppler-aware attention module and
a multi-resolution feature enhancement block that sharpened the range dimension
of ship signatures. The \gls{rd}-domain detection results provide a complementary
perspective to the optical-domain approach of the proposed system, suggesting
that future multi-modal variants could fuse optical and radar detections for
improved coverage in adverse visibility.
 
Wu et al. \cite{Wu2025} developed a two-stage target detection framework for
compact \gls{hfswr} that combined a space-to-depth
\gls{yolov8} first stage with a multi-frame Vision Transformer (\gls{vit}) second stage.
The first stage processed individual range-Doppler frames to produce candidate
targets, while the second stage aggregated temporal evidence across multiple frames
using a \gls{vit} to confirm genuine ship targets and reject clutter. Published in IEEE
Journal of Selected Topics in Applied Earth Observations and Remote Sensing,
this two-stage design is conceptually aligned with the proposed system's layered
architecture, where \gls{yolov8} provides per-frame detections that are subsequently
processed by \gls{deepocsort}'s multi-frame trajectory integrator.
 
Niu et al. \cite{Niu2026} proposed LD-\gls{yolo}, a lightweight dynamic
convolution-based \gls{yolov8}n framework for robust ship detection in \gls{sar} imagery,
published in IEEE Geoscience and Remote Sensing Letters. Dynamic convolution
replaced standard static kernels with input-conditioned kernels that adapted
their spatial weights based on the local image content, allowing a small network
to achieve the representational flexibility of a much larger static model.
The approach was specifically evaluated under varying sea-state conditions,
where wave clutter intensity varies significantly and can degrade fixed-kernel
feature extraction. The dynamic convolution concept is relevant to the proposed
system's stabilisation stage, which similarly adapts the feature reference frame
dynamically to compensate for platform motion.
 
Ai et al. \cite{Ai2025} proposed SOD-Net, a small ship object detection network
for \gls{sar} images, published in IEEE Geoscience and Remote Sensing Letters. The
network addressed the fundamental challenge that small ships occupy fewer than
0.01\% of the total image pixels, making their feature signatures easily
overwhelmed by background clutter and sidelobe artefacts. The authors introduced
a dilated residual feature pyramid combined with a super-resolution auxiliary
branch that upsampled small-target feature maps before the detection head,
providing additional spatial resolution for the regression stage. Performance
gains were concentrated on the small-ship category, confirming the value of
dedicated network components for size-specific target classes, a design principle
relevant to future enhancements of the proposed system's detection module.
 
Ou et al. \cite{Ou2025} developed DHLNet, a dynamic hierarchical lightweight
network for enhanced ship detection in remote sensing images, published in IEEE
Journal of Selected Topics in Applied Earth Observations and Remote Sensing.
The network dynamically reconfigured its hierarchical feature routing at each
forward pass based on a lightweight complexity estimator that measured the local
structural complexity of the input image. In low-complexity regions such as open
ocean, the network routed features through a shallow path for efficiency; in
high-complexity regions such as harbours and coastlines, it activated deeper paths
for precision. This dynamic depth allocation mirrors the adaptive prediction-step
scaling implemented in the proposed system's trajectory forecasting module, where
the prediction horizon is adjusted based on the measured vessel speed.
 
Liu et al. \cite{Liu2025c} addressed multi-\gls{aav}
ship detection in \gls{mec} scenarios, proposing a distributed
\gls{yolov8}-based detection framework that partitioned inference layers between
onboard \gls{aav} processors and edge servers to minimise end-to-end latency, published
in IEEE Transactions on Intelligent Transportation Systems. The system demonstrated
real-time ship localisation across multiple simultaneous aerial viewpoints, with
a consensus mechanism that resolved conflicting detections when multiple AAVs
observed the same vessel from different angles. This multi-viewpoint fusion
architecture represents a natural extension of the single-camera system proposed
in the present work, with the consensus mechanism analogous to the multi-session
tracking logic of the \gls{ais} dashboard.
 
Tang et al. \cite{Tang2025} introduced BESW-\gls{yolo}, a lightweight \gls{sar} image
detection model based on \gls{yolov8}n optimised for complex deployment scenarios,
published in IEEE Journal of Selected Topics in Applied Earth Observations and
Remote Sensing. The model incorporated a bi-level elastic stride-windowed feature
extraction module that processed \gls{sar} image chips at two stride rates simultaneously,
enabling the detection of both fine-resolution ship details and coarse-resolution
contextual layout within a single forward pass. Evaluation on the \gls{ssdd} and \gls{hrsid}
datasets demonstrated competitive accuracy with 40\% fewer parameters than
\gls{yolov8}s, confirming the efficiency of stride-windowed feature extraction. The
dual-resolution feature strategy reflects the same principle of simultaneous
coarse and fine feature representation that underlies \gls{yolov8}'s multi-scale
prediction heads used in the proposed system.
 
Zhang et al. \cite{Zhang2026} presented SegS-\gls{yolo}, a \gls{yolo}-based instance
segmentation approach for small-scale ship targets in \gls{sar} images, published
in IEEE Journal of Selected Topics in Applied Earth Observations and Remote
Sensing. The segmentation masks provided more precise ship boundary estimates
than bounding boxes, enabling accurate computation of ship length and beam from
the mask geometry. The authors demonstrated that mask-derived size estimates
correlated with ground-truth ship dimensions from \gls{ais} records, establishing
a promising link between vision-based detection output and \gls{ais}-reported vessel
characteristics. This \gls{ais}-vision cross-validation methodology directly motivates
the integration of \gls{ais} data with the computer-vision detection results in the
proposed system's unified dashboard.
 
\subsubsection{\textbf{Cross-Domain and Application-Specific Deep Learning Detection}}
Caricchio et al. \cite{Caricchio2025} applied the \gls{yolov8} neural network to
non-collaborative vessel detection using Sentinel-1 \gls{sar} data, published in IEEE
Geoscience and Remote Sensing Letters. Non-collaborative vessels --- those that
have deactivated or manipulated their \gls{ais} transponders --- represent a critical
surveillance target that radar-based systems often fail to identify definitively.
The study demonstrated that \gls{yolov8} could detect such vessels directly from \gls{sar}
amplitude imagery without any \gls{ais} assistance, achieving detection rates comparable
to specialised matched-filter-based radar detectors. The complementarity between
vision-based detection and \gls{ais} data fusion demonstrated in this paper directly
motivates the proposed system's dual-stream architecture, where \gls{ais} data enriches
but does not replace the computer-vision detection pipeline.
 
Cheng et al. \cite{Cheng2025} addressed automatic shipping container attribute
spotting, developing a fast and lightweight system for recognising container
identifiers and size codes from port surveillance imagery, published in IEEE
Transactions on Instrumentation and Measurement. While focused on containers
rather than vessels, the work characterised a production deployment of \gls{yolov8}
in a maritime logistics context, demonstrating the framework's versatility across
different maritime target categories. The deployment insights regarding illumination
handling under the mixed artificial and natural lighting conditions of port
environments are directly applicable to the proposed surveillance system's
optical detection module.
 
Xu et al. \cite{Xu2025a} proposed a multistage fusion object detection system
combining marine radar and camera inputs for complex maritime contexts, published
in IEEE Transactions on Instrumentation and Measurement. The fusion pipeline
aligned radar range-bearing detections with camera image regions using a calibrated
extrinsic transformation, then fused confidence scores from both modalities using
a Dempster-Shafer evidence combination rule. The approach achieved higher detection
completeness than either modality alone, with radar providing coverage in visual
obscuration conditions and the camera providing classification capability in clear
conditions. The complementary strengths of radar and camera modalities highlighted
in this study directly motivate the proposed system's future integration direction,
where the vision pipeline could be augmented with radar detection inputs.
 
Nambiar et al. \cite{Nambiar2022} explored efficient ship detection using deep
learning techniques applied to both \gls{sar} images and lateral-view ship photographs,
presenting one of the earlier systematic comparisons of \gls{yolo} variants for maritime
detection. The study benchmarked YOLOv4 and YOLOv5 variants and found that the
lateral-view optical modality generally yielded higher detection precision while
\gls{sar} provided superior coverage under foggy or night conditions. The finding that
no single modality dominates across all conditions provided early evidence for
the multi-modal fusion direction subsequently pursued in the literature. This
foundational comparison contextualises the present system's choice of optical
video as the primary detection modality, appropriate for clear daytime operation,
with \gls{ais} as a complementary data layer.
 
Geng et al. \cite{Geng2023} presented a vision-based ship speed measurement
method using deep learning, specifically estimating vessel speed from monocular
camera footage by combining object detection with optical flow-based displacement
estimation. The deep learning model produced calibrated speed estimates with a
mean absolute error of less than 0.5 knots compared to \gls{ais}-reported speed over
ground. This result is particularly relevant to the proposed system, which
similarly derives speed estimates from inter-frame centroid displacements, and
provides external validation that vision-based speed estimation can approach
\gls{ais}-level accuracy with appropriate calibration.
 
Ranasinghe et al. \cite{Ranasinghe2023} proposed integrating low-cost acoustic
sensors with visual detection for autonomous vessel detection, fusing sonar-derived
range information with camera-based appearance features. The sensor fusion strategy
improved detection in conditions where one modality was degraded, with acoustic
data providing bearing constraints that reduced the camera-only detection search
space. The broader theme of complementary sensing modalities illustrated in this
paper reinforces the design rationale of the proposed system, which pairs the
optical detection pipeline with \gls{ais} data to provide context that single-modality
approaches cannot offer.
 
Zhang et al.\ \cite{Zhang2025a} developed an autonomous landing method for
ship-borne \glspl{uav} based on visual recognition of landing
pad markers, published in IEEE Access. While not a ship detection paper in the
conventional sense, the study demonstrated the precision and reliability of
\gls{yolov8}-based visual localisation in the maritime context, where the \gls{uav}'s camera
experienced significant motion from ship heave and roll. The successful handling
of camera motion in this application is directly analogous to the challenge
addressed by the drone stabiliser in the proposed system, validating the
Lucas-Kanade optical-flow approach for maritime platform motion compensation.
 
Chen and Wang \cite{Chen2024a} proposed a shuffled grouping cross-channel
attention-based bilateral-filter-interpolation deformable ConvNet for benthonic
organism detection, published in IEEE Transactions on Artificial Intelligence.
The underwater detection context shares several challenges with maritime surface
surveillance, including irregular target shapes, cluttered backgrounds, and
variable illumination from water-surface reflections. The authors' bilateral-filter
preprocessing stage, which smoothed background texture while preserving target
boundaries, is conceptually analogous to the stabilisation and normalisation
preprocessing stage in the proposed system, reinforcing the general principle
that domain-appropriate preprocessing significantly improves downstream detection
performance.
 
Dasaramoole Prakash et al. \cite{DasaramoolePrakash2024} presented a lightweight
and efficient \gls{yolov8} with a residual attention mechanism for medical image
classification, published in IEEE Access, demonstrating the transferability of
the \gls{yolov8} architecture and residual attention blocks across vastly different
application domains. The study's core finding that residual attention connectivity
consistently improved classification confidence calibration, regardless of the
target domain, provides domain-agnostic justification for the attention module
augmentations applied to \gls{yolov8} in the maritime detection literature surveyed
above, reinforcing their architectural validity beyond the maritime context.
 
Cai et al. \cite{Cai2026} developed a contactless intelligent anti-interference
lung nodule detection method for early disease detection, published in IEEE
Journal of Biomedical and Health Informatics. While the application is medical,
the paper made a methodological contribution of direct relevance to the proposed
system: the use of a multi-scale feature aggregation strategy that maintained
high detection sensitivity for very small targets in cluttered backgrounds.
The demonstrated efficacy of multi-scale aggregation for small target detection
in high-clutter environments directly parallels the small vessel detection
challenge in maritime surveillance and reinforces the architectural design
decisions underlying the \gls{yolov8} feature pyramid used in the proposed system.
 
Zhai et al. \cite{Zhai2025} developed DBFFNet, a dual-branch feature fusion
framework for high-precision flame and smoke detection in early fire scenarios,
published in IEEE Access. The study demonstrated a sophisticated feature fusion
design principle whereby spatial and channel attention branches were trained
independently and fused at the decision level, achieving superior detection of
diffuse, irregular targets against cluttered backgrounds. The dual-branch attention
architecture provides design inspiration for future improvements to the proposed
system's detection backbone for handling diffuse maritime targets such as oil spills
or low-contrast wakes, where separate spatial and channel attention pathways could
independently specialise their feature extraction strategies.
 
\subsection{Multi-Object Tracking and \texorpdfstring{\gls{reid}}{ReID}}
The multi-object tracking problem in maritime surveillance extends beyond
detection, requiring the maintenance of consistent vessel identities across
time, the resolution of occlusion events, and the management of identity
switches in cluttered environments. The following papers trace the evolution
of Simple Online and Real-time Tracking (\gls{sort})-based methods toward
appearance-augmented trackers and their application to maritime and related
surveillance scenarios.
 
\subsubsection{\textbf{\gls{sort}-Based Tracking Frameworks and Extensions}}
Belmouhcine et al. \cite{Belmouhcine2021} introduced Robust Deep \gls{sort}, a
strengthened variant of the \gls{deepsort} tracking framework that replaced the standard
Euclidean distance matching metric with a learned robust distance that downweighted
outlier detections produced by partial occlusion. The authors demonstrated that
conventional \gls{deepsort} was brittle to detection noise, frequently generating
erroneous tracklet terminations when a vessel was momentarily obscured. The robust
distance metric, trained on a curated set of occlusion examples, reduced identity
switch events by over 30\% on the MOT17 benchmark. This work is directly
foundational to the \gls{deepocsort} tracker employed in the proposed system, which
inherits the appearance-matching formulation and extends it with an online
clustering update for the \gls{reid} feature database.
 
Belmouhcine et al. \cite{Belmouhcine2022} further extended their tracking
framework in Fully Deep \gls{sort}, published at the International Conference on Pattern
Recognition. The re-identification component was upgraded to use attention-based
similarity learning rather than explicit metric learning, producing a more efficient
and accurate matching stage by eliminating the hand-crafted triplet loss training
pipeline. The attention-based \gls{reid} module computed pairwise similarity scores
directly from encoded appearance features without requiring pre-defined distance
functions, achieving superior recall at lower false-positive rates on the DukeMTMC
and Market-1501 benchmarks. This progressive shift from hand-crafted to learned
matching components in tracking frameworks informs the design of the proposed
system, which retains the Kalman filter for \gls{gps}-space smoothing but uses the
learned \gls{osnet} \gls{reid} for appearance matching.
 
Wang et al. \cite{Wang2025a} proposed an efficient maritime traffic flow analysis
system combining \gls{yolov8} detection with \gls{deepsort} tracking, published at IAEAC 2025.
The system was deployed on a coastal camera network covering a busy shipping lane
and demonstrated continuous vessel tracking at 25 fps across a 1.2-kilometre
observable corridor. Vessel count statistics aggregated from tracking results were
used to model traffic flow patterns, with lane occupancy and speed distributions
extracted from historical track data. The work established a direct precedent for
the proposed system's deployment scenario and confirmed that \gls{yolov8} plus
\gls{deepsort}-family trackers can sustain real-time operation in high-traffic maritime
environments, validating the technical approach at a systems level.
 
\subsubsection{\textbf{Maritime-Specific Tracking Systems}}
Muslim et al. \cite{Muslim2025} developed a system for automatic vessel detection
and tracking in the Port of Split, integrating a \gls{yolov8} detector with a geometric
centroid tracker adapted for the constrained geometry of a port basin, published
at SpliTech 2025. The port-specific tracker incorporated prior knowledge of vessel
berth locations and approach corridors to reject detections inconsistent with
plausible vessel trajectories, reducing false-positive tracks generated by wave
reflections and floating debris. The system operated continuously over a 72-hour
evaluation period with manually verified tracking accuracy. The principle of
incorporating domain-specific spatial priors into the tracking logic is relevant
to the proposed system, where the land-mask check in the geo-mapper similarly
enforces the physical constraint that vessels cannot occupy land areas.
 
Liu et al. \cite{Liu2024} investigated multiple ship tracking across multiple \gls{leo}
satellites, addressing the problem of inter-pass ship identity association when
the same vessel is observed by different satellite passes separated by hours,
published at the IEEE IGARSS 2024 symposium. The authors used \gls{sar}-derived ship
signatures combined with \gls{ais} records to build a cross-pass matching database,
then applied a bipartite graph matching algorithm to associate detections across
passes. The inter-pass identity association problem is conceptually analogous to
the re-identification problem addressed by \gls{osnet} \gls{reid} in the proposed system,
where a vessel temporarily exits and re-enters the camera field of view.
 
Nithya et al.\ \cite{Nithya2023} examined multi small object detection and
prioritised tracking for Navy operations using deep learning techniques, addressing
the challenge of simultaneously tracking numerous small high-speed surface targets
under engagement-decision time pressure. The prioritised tracking framework ranked
tracked objects by threat-relevant attributes including speed, bearing change,
and proximity to the protected asset, and allocated higher computational resources
to high-priority tracks. While the military application context differs from the
civilian surveillance focus of the proposed system, the principle of differentiating
tracking resources by behavioural criteria is directly applicable to future
extensions of the proposed system targeting anomalous vessel detection.
 
\subsubsection{\textbf{Appearance-Based Tracking in Related Surveillance Domains}}
Zang et al. \cite{Zang2024} applied a \gls{yolov8}-plus-\gls{deepsort} pipeline to smoke
and fire target detection and tracking in surveillance video, published at YAC
2024. Although the application domain differs from maritime surveillance, the
study provided a careful quantitative analysis of the tracking pipeline's
behaviour under appearance ambiguity conditions, specifically when multiple fire
sources produced similar colour and texture signatures causing identity confusion
in the \gls{reid} matching stage. The authors' proposed solution, a spatially constrained
matching window that prioritised proximity over appearance similarity for ambiguous
cases, is directly applicable to the vessel tracking context where closely spaced
vessels with similar hull colours can similarly confuse appearance-based matchers.
 
Wang and Kong \cite{Wang2024} developed an efficient association multiple-object
tracker for marine benthonic organism counting, addressing the challenge of tracking
slow-moving, irregularly shaped organisms in underwater video with frequent partial
occlusion. The tracker combined an intersection-over-union (IoU)-based primary
association with a secondary appearance-based re-entry mechanism for lost tracklets,
achieving high counting accuracy over long video sequences. Published at the Chinese
Control Conference, this work demonstrated that the \gls{sort}-family tracking paradigm
extends effectively to non-vessel maritime targets, reinforcing the generality of
the proposed system's tracking architecture for diverse maritime surveillance
applications.
 
Park et al. \cite{Park2025} proposed a real-time snowfall recognition and object
tracking system for autonomous driving in adverse snowy weather, published in
IEEE Access. The study's core contribution was a weather-condition-aware tracking
threshold that adjusted the \gls{reid} matching confidence threshold based on estimated
snowfall intensity, accepting weaker appearance matches in heavy snowfall where
visual degradation reduced discriminability. The principle of environment-adaptive
matching thresholds has direct applicability to maritime surveillance, where fog,
glare, and sea spray similarly degrade visual features, suggesting that a
sea-state-adaptive \gls{reid} threshold could improve the proposed system's performance
in adverse conditions.
 
\subsubsection{\textbf{Collision Avoidance and Multi-Source Maritime Tracking}}
Zhao et al. \cite{Zhao2024} studied ship-bridge collision prevention through early
warning based on multi-source data fusion, integrating \gls{ais}, radar, and camera-based
tracking data within a unified risk assessment framework, published at ACAIT 2024.
The authors demonstrated that camera-based tracking provided position updates at a
finer temporal resolution than \gls{ais}, which transmits at 2--10 second intervals for
vessels at speed, enabling faster collision risk estimation. The proposed fusion
weighting scheme gave camera-tracking data priority for short-term manoeuvre
prediction while using \gls{ais} data for vessel identity labelling and long-term trajectory
context. This temporal complementarity between camera tracking and \gls{ais} data is
directly implemented in the proposed system, where the vision pipeline provides
high-rate centroid updates and \gls{ais} provides identity and context enrichment.
 
Wan et al. \cite{Wan2025} investigated methods for multi-source sensing fusion
in dynamic maritime environments, published at ICTIS 2025. The study benchmarked
four fusion architectures --- early fusion, late fusion, cross-modal attention,
and evidence-based fusion --- for combining radar, \gls{ais}, optical camera, and
infrared sensor data in a vessel situation awareness system. Cross-modal attention
fusion consistently outperformed the other strategies, particularly in scenarios
where one sensor modality was temporarily degraded. The benchmark results provide
a systematic guidance framework for future extensions of the proposed system from
its current dual-stream vision-plus-\gls{ais} architecture toward full multi-sensor
fusion.
 
Suliva et al. \cite{Suliva2022} presented a system for classification and counting
of ships using the YOLOv5 algorithm, published at ICCIS 2022, establishing an
early baseline for vision-based vessel enumeration in maritime surveillance. The
study provided detailed analysis of detection confusion between vessel categories
and recommended domain-specific fine-tuning of the detection backbone on
vessel-type-labelled maritime datasets. The classification confusion analysis
directly informs the class mapping employed in the proposed system, which defaults
to \gls{coco} class index 8 (boat) as a coarse but reliable vessel-presence indicator
rather than attempting fine-grained type classification without domain-specific
training data.
 
Zhu and Zhang \cite{Zhu2024} applied YOLOv5 to multi-target tracking and detection
on sea surfaces, published at ICIPMC 2024. The study investigated how wave motion
and sun glare affected bounding box stability across frames, causing jittery track
histories even for stationary vessels. The authors applied a temporal median filter
to stabilise bounding box coordinates before passing them to the tracking stage,
achieving smoother track histories and improved speed estimation. This temporal
smoothing approach is complementary to the spatial smoothing implemented in the
proposed system's centroid history function, which applies a moving-average filter
to centroid positions before heading and speed computation.
 
\subsection{Vessel Trajectory Prediction and Path Forecasting}
Accurate prediction of future vessel positions is a critical capability for
proactive maritime traffic management, collision avoidance, and anomaly detection.
Research in this domain has progressed from classical kinematic models through
recurrent neural network-based sequence learning to recent transformer and
multimodal approaches. The following papers are reviewed within thematic groups
covering \gls{lstm}-based, transformer-based, and hybrid statistical approaches.
 
\subsubsection{\textbf{Recurrent Neural Networks for Trajectory Forecasting}}
Capobianco et al. \cite{Capobianco2021} provided one of the seminal studies of
deep learning methods for vessel trajectory prediction based on Recurrent Neural
Networks (\glspl{rnn}), published in IEEE Transactions on Aerospace and Electronic Systems.
The authors systematically compared vanilla \gls{rnn}, \gls{lstm}, and Gated Recurrent Unit
(\gls{gru}) architectures for predicting vessel position sequences derived from \gls{ais}
records, with a focus on long-horizon prediction accuracy up to 60 minutes ahead.
\gls{lstm} consistently outperformed the other architectures across all prediction horizons
and traffic densities, attributed to its gating mechanism's ability to selectively
retain relevant long-term trajectory context while discarding transient manoeuvre
artefacts. This foundational result directly motivates the two-layer \gls{lstm} architecture
adopted in the proposed system's trajectory forecasting module.
 
Yang et al. \cite{Yang2022} proposed \gls{ais}-based intelligent vessel trajectory
prediction using \gls{bilstm}, published in IEEE Access. The
\gls{bilstm} processed \gls{ais} position sequences in both forward and backward temporal
directions, allowing the model to incorporate contextual information from both
ends of the observed trajectory during training, producing richer internal
representations of trajectory shape than unidirectional \glspl{lstm}. Experiments on
the South China Sea \gls{ais} dataset demonstrated that \gls{bilstm} predictions achieved
a 12\% lower \gls{mae} at 30-minute prediction horizons compared
to standard \gls{lstm}. While the proposed system employs a unidirectional \gls{lstm} for
online prediction where future data is unavailable, the \gls{bilstm} results establish
an upper-bound reference for the expected accuracy improvement from incorporating
longer historical context.
 
Chondrodima et al. \cite{Chondrodima2023} developed an efficient \gls{lstm} neural
network-based framework for vessel location forecasting, published in IEEE
Transactions on Intelligent Transportation Systems. The key contribution was a
data preprocessing pipeline that transformed raw \gls{ais} latitude and longitude
sequences into a normalised displacement representation, effectively decoupling
the prediction problem from the geographic coordinate system and enabling a single
trained model to generalise across different ocean regions without retraining.
The displacement normalisation strategy is directly implemented in the proposed
system's \gls{lstm} input pipeline, where \gls{gps} position inputs are converted to velocity
tuples before being fed to the network to ensure geographic generalisation.
 
\subsubsection{\textbf{Transformer and Context-Aware Prediction Approaches}}
Nguyen and Fablet \cite{Nguyen2024} proposed a transformer network with sparse
augmented data representation and cross-entropy loss for \gls{ais}-based vessel trajectory
prediction, published in IEEE Access. The transformer architecture replaced the
\gls{lstm}'s sequential hidden state with a parallel self-attention mechanism that could
selectively focus on any past timestep regardless of temporal distance, overcoming
the gradient vanishing limitation of \gls{lstm} for very long sequences. The cross-entropy
loss formulation discretised the prediction space into a geographic grid, enabling
probabilistic uncertainty estimates alongside point predictions. The probabilistic
output modality represents a natural and important extension for future versions
of the proposed system, enabling uncertainty-aware collision avoidance alerting.
 
Zhang et al. \cite{Zhang2024} developed a dynamic context-aware approach for
vessel trajectory prediction based on multi-stage deep learning, published in
IEEE Transactions on Intelligent Vehicles. The multi-stage pipeline first classified
the vessel's operational mode --- cruising, manoeuvring, or anchoring --- from the
recent trajectory segment, then selected a specialised prediction sub-model trained
for that operational mode. The context-aware switching mechanism produced significantly lower prediction errors during mode transitions compared to a single global model, because the sub-models were not required to simultaneously represent the disparate dynamics of cruising and manoeuvring. The operational mode classification concept is directly applicable to the proposed system's trajectory module, where kinematic prediction is selected during the cold-start phase before \gls{lstm} predictions become reliable, analogous to an anchoring or slow-manoeuvre operational mode.
 
Zhang et al. \cite{Zhang2026a} presented a unified multimodal vessel trajectory
prediction framework with explainable navigation intention, published in IEEE
Transactions on Intelligent Transportation Systems. The framework jointly modelled
\gls{ais} kinematics, vessel type attributes, and environmental factors such as weather
and traffic density within a single transformer encoder, producing intention-level
explanations alongside trajectory predictions that described the vessel's likely
navigational objective. The explainability dimension addresses a critical requirement
for operational deployment of AI-based surveillance systems, where operators must
be able to assess prediction rationale. The intention modelling concept represents
an important future extension for the proposed system's trajectory prediction module.
 
Han et al. \cite{Han2024} addressed interaction-aware short-term marine vessel
trajectory prediction with deep generative models, published in IEEE Transactions
on Industrial Informatics. The core contribution was a \gls{cvae} that modelled the joint distribution of future positions for
all vessels in a scene simultaneously, capturing inter-vessel interaction effects
such as overtaking, channel priority, and collision avoidance manoeuvres that
single-vessel predictors systematically ignore. The \gls{cvae} approach generated a
distribution of plausible future trajectories rather than a single deterministic
prediction, enabling risk assessment based on the overlap of predicted trajectory
distributions for pairs of vessels. The interaction modelling capability is absent
from the proposed system's current \gls{lstm} module, where each vessel's trajectory is
predicted independently, and represents the most significant avenue for future
enhancement.
 
\subsubsection{\textbf{\gls{ais}-Based Statistical and Hybrid Forecasting Methods}}
Zhang et al. \cite{Zhang2022} provided a comprehensive survey of vessel trajectory
prediction approaches in maritime transportation, published in IEEE Transactions on
Intelligent Transportation Systems. The survey covered statistical methods including
Gaussian Processes and polynomial regression, machine learning methods including
Support Vector Regression (\gls{svr}) and Random Forests, and deep learning methods
including \gls{lstm}, \gls{gru}, and Transformers, providing a unified evaluation framework
based on prediction horizon, data density, and generalisation to unseen routes.
The survey identified weighted polynomial regression as consistently competitive
for short horizons of less than five minutes despite its simplicity, directly
validating the use of weighted polynomial regression as the kinematic fallback
predictor in the proposed system, and identified \gls{lstm} as the dominant approach for
medium horizons of five to thirty minutes, validating the primary \gls{lstm} predictor.
 
Zygouras et al. \cite{Zygouras2024} introduced EnvClus*, an algorithm for
extracting common navigational pathways from historical \gls{ais} data to support
effective vessel trajectory forecasting, published in IEEE Access. The algorithm
clustered \gls{ais} trajectories by their environmental context such as fairways,
anchorage zones, and fishing grounds, and computed representative pathway templates
for each cluster. Trajectory prediction was performed by identifying the most
likely pathway template for the current vessel and extrapolating along it, producing
route-conformant predictions that outperformed purely data-driven approaches on
structured navigational routes. The pathway template concept is directly applicable
to future extensions of the proposed system in geographic areas with well-defined
shipping lanes.
 
Liu et al. \cite{Liu2022} presented deep learning-powered vessel trajectory
prediction for improving smart traffic services in the \gls{miot}, published in IEEE Transactions on Network Science and Engineering. An
attention-augmented \gls{lstm} was proposed that weighted historical \gls{ais} observations
by their estimated reliability based on timestamp regularity and reported accuracy,
producing more accurate predictions from sparse or irregular \gls{ais} streams. The
reliability-weighted input mechanism addresses the same data quality challenge
that motivates the \gls{gps}-space Kalman filter in the proposed system, which smooths
noisy pixel-to-\gls{gps} projections before they enter the \gls{lstm} input pipeline.
 
Chen et al. \cite{Chen2025} addressed machine learning-based reconstruction of
missing vessel trajectories, published in IEEE Transactions on Vehicular Technology.
\gls{ais} data streams frequently contain gaps caused by signal loss, congestion, or
deliberate transponder deactivation. The authors trained an encoder-decoder \gls{lstm}
that could interpolate missing trajectory segments given the available surrounding
context, recovering trajectory continuity with sub-kilometre positional error.
The trajectory reconstruction capability is directly analogous to the kinematic
fallback predictor in the proposed system, which maintains position estimates for
vessels during temporary occlusion events where the camera-based tracker produces
no detections.
 
Aouedi et al. \cite{Aouedi2026} developed an \gls{ais}-based hybrid vessel trajectory
prediction framework for enhanced maritime navigation, published in IEEE Transactions
on Machine Learning in Communications and Networking. The hybrid architecture combined
a physics-informed kinematic model with a data-driven \gls{lstm}, fusing their predictions
through a learnable weighting network that dynamically allocated prediction authority
based on the estimated reliability of each component. During open-ocean straight-line
segments, the kinematic model dominated; during port approach manoeuvres with complex
trajectory curvature, the \gls{lstm} dominated. The hybrid fusion approach achieved lower
prediction error than either component in isolation across all route segments. This
hybrid design principle is directly embodied in the proposed system, where kinematic
polynomial prediction and \gls{lstm} prediction are combined with a cold-start switching
rule.
 
\subsection{\texorpdfstring{\gls{ais}}{AIS} Data Integration and Maritime Intelligence}
The \gls{ais} remains the primary source of authoritative
vessel identity and kinematic information in maritime traffic management. The
following papers examine how \gls{ais} data quality, geolocation accuracy, and
multi-source fusion contribute to the broader maritime surveillance mission,
and how these findings directly shaped the design of the proposed system's
\gls{ais} dashboard component.
 
\subsubsection{\textbf{\gls{ais} Data Quality and Geolocation Correction}}
Wu et al. \cite{Wu2024} developed an \gls{ais} data-guided geolocation correction method
for Low-Orbit Satellite Remote Sensing imagery, published in IEEE Journal of Selected
Topics in Applied Earth Observations and Remote Sensing. The method used \gls{ais}-reported
position records to serve as ground truth reference points for correcting the
geolocation errors inherent in rapidly moving satellite platforms, where attitude
estimation errors and orbital parameter uncertainties produce systematic spatial
offsets in the imagery. The \gls{ais}-guided correction achieved sub-pixel geolocation
accuracy on the SENTINEL-2 dataset, enabling reliable fusion of satellite imagery
and \gls{ais} records. This geolocation correction methodology is conceptually aligned
with the homography-based geo-mapper in the proposed system, which similarly uses
known reference coordinates to establish a transformation between pixel space and
geographic space.
 
Guan et al. \cite{Guan2023} studied the recognition of multiple Maritime Mobile
Service Identity (\gls{mmsi}) codes on a single vessel using \gls{ais} data, published in
IEEE Access. \gls{mmsi} code manipulation, where vessels transmit multiple or falsified
identity codes to evade surveillance, is a significant challenge for \gls{ais}-based
monitoring. The authors proposed a machine learning classifier that identified
anomalous \gls{mmsi} transmission patterns by analysing the consistency of reported
kinematic parameters --- speed, heading, and position rate of change --- across
transmissions from the same vessel. The identity anomaly detection methodology
is directly relevant to the proposed system's \gls{ais} dashboard, where \gls{mmsi}-based
vessel identification could be augmented with behaviour-based consistency checks
to flag potential identity spoofing events.
 
\subsubsection{\textbf{Multi-Vessel Data Fusion and Graph-Based Association}}
Lu et al. \cite{Lu2026} proposed a graph learning-driven multi-vessel association
framework fusing multimodal data for maritime intelligence, published in IEEE
Transactions on Intelligent Transportation Systems. The framework modelled all
vessels in a geographic region as nodes in a dynamic graph, with edges representing
spatial proximity and potential interaction relationships. Graph neural network
message passing propagated identity and trajectory information across the graph,
enabling the association of the same vessel detected by different sensors even
when the sensor-specific identifiers did not match directly. The graph-based
multi-source association is a sophisticated extension of the \gls{ais}-plus-vision fusion
implemented in the proposed system's unified dashboard, representing a natural next
step for integrating additional sensor modalities.
 
Wan et al. \cite{Wan2025} provided a systematic benchmark of multi-source sensing
fusion methods for dynamic maritime environments, covering radar, \gls{ais}, optical, and
infrared modalities. The study quantified the performance of four fusion architectures
across twelve representative maritime scenarios ranging from clear open-ocean
conditions to foggy port approaches. Cross-modal attention fusion consistently
outperformed the other strategies, particularly in scenarios where one sensor
modality was temporarily degraded. The standardised metrics proposed in this study
provide an evaluation framework applicable to the proposed system's unified
dashboard, enabling principled comparison against future multi-sensor variants.
 
Zhao et al. \cite{Zhao2024} developed a multi-source data fusion approach for
ship-bridge collision prevention early warning, combining \gls{ais} trajectory forecasts
with camera-based vessel proximity estimates. The system computed a dynamic risk
index incorporating predicted closest point of approach, vessel speed, and turning
rate, triggering alerts when the index exceeded a configurable threshold. Field
trials demonstrated reliable early warning with a median lead time of 4.2 minutes
before estimated collision, sufficient for navigational intervention. The collision
avoidance application represents one of the most operationally significant downstream
uses of the proposed system's trajectory prediction capability.

\subsubsection{\textbf{\gls{ais} Data Filtering, Outlier Detection, and Sensor Fusion}}

Urakami et al. \cite{Urakami2018} presented a study on location information
screening methods for ship applications using \gls{ais} recorded data, published at the
International Conference on Broadband Communications for Next Generation Networks
and Multimedia Applications (CoBCom). The authors identified that raw \gls{ais} position
logs frequently contain erroneous position reports arising from \gls{gps} multi-path
interference, antenna misalignment, and transponder firmware faults, producing
sudden large positional jumps that corrupt trajectory analysis and downstream
prediction models. A rule-based screening pipeline was proposed that flagged
position records whose implied speed-over-ground or change-in-heading exceeded
physically plausible thresholds for the vessel category, discarding outlier records
and interpolating replacement positions from the surrounding valid sequence. The
study demonstrated that applying the screening stage upstream of any trajectory
analysis reduced the variance of derived kinematic estimates by a significant
margin. This data cleaning rationale maps directly onto the land-mask outlier
rejection and \gls{gps}-space Kalman filter stages of the proposed system, which together
serve the same function for vision-derived \gls{gps} positions: discarding physically
implausible coordinate projections before they corrupt the \gls{lstm} input sequence
and trajectory prediction output.

Lian et al.\cite{Lian2019} proposed a ship \gls{ais} trajectory estimation method based on the Particle Filter algorithm, published at the International Conference on Intelligent Human-Machine Systems and Cybernetics (IHMSC). The core motivation was the estimation of vessel position during \gls{ais} data gaps, which arise due to \gls{vhf} radio coverage holes, signal congestion in high-traffic areas, and deliberate transponder deactivation. The particle filter maintained a probability distribution over possible vessel states by propagating a set of weighted particles forward in time using a kinematic motion model, then reweighting them when an \gls{ais} observation became available. This Bayesian sequential estimation approach outperformed linear interpolation and Kalman-only methods during prolonged gap intervals, particularly when vessels executed course changes during the gap period that violated the constant-velocity assumption. The relevance to the proposed system is direct and significant: the \gls{gps}-space Kalman filter employed in the geo-mapping stage is a linear special case of the general Bayesian state estimator demonstrated by this paper, and the particle filter approach represents a natural upgrade path for handling non-linear vessel manoeuvres that the Kalman filter's constant-velocity model approximates imperfectly.

Wang and Chen \cite{Wang2024a} developed a ship \gls{ais} trajectory outlier identification model under high-speed sampling conditions, published at the International Symposium on Computer Technology and Information Science (ISCTIS). High-speed \gls{ais} sampling, where position records are generated at sub-second intervals rather than the standard 2--10 second intervals, produces dense trajectory datasets that amplify the presence of \gls{gps} jitter and transient measurement noise. The authors introduced a density-based spatial clustering algorithm adapted for the time-ordered structure of \gls{ais} sequences, which identified outlier records as those whose position deviated from the local trajectory manifold by more than a learned threshold. The model was evaluated on high-frequency \gls{ais} recordings from a congested port approach channel, demonstrating consistent outlier rejection with a low false-positive rate on genuine manoeuvres. The dense sampling scenario addressed in this paper is directly analogous to the frame-rate \gls{gps} position stream produced by the proposed system's geo-mapper, where each video frame generates a new \gls{gps} projection that exhibits the same jitter and noise characteristics as high-frequency \gls{ais} data. The outlier identification strategy described by Wang and Chen complements the Kalman filter smoothing stage of the proposed system and suggests a clustering-based pre-screening step as a candidate enhancement for future versions.

Chen \cite{Chen2023a} addressed the efficient track association of ships through \gls{ais} and shipboard video fusion via ad hoc networks, published at the International Conference on Computer and Communications (ICCC). The study tackled the fundamental challenge of associating vision-tracked vessel bounding boxes with \gls{ais}-reported vessel identities when the two data streams differ in coordinate frame, update rate, and identity representation. The proposed fusion architecture transmitted \gls{ais} records and compressed video features across a shipboard ad hoc network with limited bandwidth, then performed association on a shore-side server using a bipartite matching formulation that minimised a joint cost combining predicted position deviation and appearance feature distance. The system demonstrated reliable track association in a multi-vessel encounter scenario where three vessels were simultaneously visible in the video stream and broadcasting \gls{ais} signals. This \gls{ais}-to-video association problem is precisely the integration challenge addressed by the proposed system's unified dashboard, where vision-pipeline track identifiers
and \gls{ais} \gls{mmsi} codes must be linked to provide operators with a single coherent vessel picture. The bipartite matching formulation
demonstrated by Chen provides a principled algorithmic foundation for implementing this association in future versions of the proposed dashboard.

Gurgel et al.\ \cite{Gurgel2010} evaluated a \gls{hfswr} ship detection and tracking algorithm by systematic comparison against \gls{ais} records and \gls{sar} imagery, published at the International WaterSide Security Conference. The study demonstrated that \gls{ais} data served as an authoritative ground truth reference for validating radar-based vessel detections, enabling the identification of both missed detections (vessels present in \gls{ais} but absent from the radar track) and false alarms (radar tracks with no corresponding \gls{ais} record). The \gls{sar} imagery provided complementary geometric validation of vessel positions, resolving cases where \gls{ais}-reported coordinates disagreed with radar-estimated positions due to \gls{ais} propagation delay or \gls{gps} error. The tri-modal comparison framework --- radar tracks, \gls{ais} records, and \gls{sar} imagery --- established a benchmark methodology for evaluating heterogeneous maritime surveillance systems that remains directly applicable to validating the proposed system. In particular, the paper's use of \gls{ais} as a validation ground truth for vision-based detection, rather than as a fusion input, represents the most rigorous evaluation strategy available for the proposed system in the absence of purpose-built maritime video ground truth datasets, and directly motivates the use of \gls{ais} position records as a reference against which the geo-mapper's \gls{gps} projections can be calibrated and validated in future experimental evaluation of the pipeline. 

\setlength{\parskip}{1em}


\subsection{Summary of Literature Review and Identified Research Gaps}

The foregoing review reveals several consistent themes and residual gaps in the
existing literature. In ship detection, the \gls{yolov8} family has emerged as the
dominant framework across optical, \gls{sar}, and infrared modalities
\cite{Ning2024, Gao2024, Gong2024, Chen2024, Yu2023, Liu2025, Fu2025,
Zhou2024, Yuan2024, Shi2025, Chen2023, Cao2024, Deshmukh2025, Caricchio2025},
with a consistent finding that architectural modifications to the backbone, neck,
and head yield incremental but reliable accuracy improvements. However, the
majority of these studies address still-image or \gls{sar} detection rather than
continuous optical video surveillance from a moving drone platform, where camera
motion compensation is an essential prerequisite to stable detection performance.

In multi-object tracking, the literature has validated \gls{sort}-based appearance
trackers for maritime scenarios \cite{Wang2025a, Muslim2025,
Belmouhcine2021, Belmouhcine2022}, but few works address the specific challenge
of tracking vessels through the combined perturbations of drone camera motion,
water surface reflections, and mutual occlusion in dense traffic. The absence
of a drone-stabilised, appearance-tracking pipeline for maritime video surveillance
represents a clear gap directly addressed by the proposed system.

In trajectory prediction, \gls{lstm}-based approaches have been validated extensively
on \gls{ais} data \cite{Capobianco2021, Yang2022, Chondrodima2023, Liu2022, Aouedi2026},
but the transition from \gls{ais}-derived \gls{gps} positions to vision-derived pixel centroids
projected into \gls{gps} space introduces coordinate noise that none of the reviewed
works addresses with a Kalman filter smoothing stage prior to \gls{lstm} input.
Similarly, the kinematic-to-\gls{lstm} cold-start management strategy employed in the
proposed system is absent from reviewed prediction frameworks
\cite{Zhang2022, Han2024, Nguyen2024, Zhang2024}, which uniformly assume a
minimum history length is always available from the outset of tracking.

In \gls{ais} data quality, reviewed works have identified the persistent problem of
erroneous, missing, and high-frequency-noise-corrupted position records in raw
\gls{ais} streams \cite{Urakami2018, Lian2019, Wang2024a}. Urakami et al.\
\cite{Urakami2018} demonstrated that kinematic threshold-based screening is
necessary before any trajectory analysis, Lian et al.\ \cite{Lian2019} showed
that particle filter estimation reliably bridges \gls{ais} data gaps during coverage
holes and transponder outages, and Wang and Chen \cite{Wang2024a} established
that density-based outlier rejection is essential under high-frequency sampling
conditions. Despite these findings, none of the reviewed works applies equivalent
data cleaning to the noisier coordinate stream produced when pixel-space vessel
centroids are projected into \gls{gps} space through a camera model --- precisely the
scenario introduced by the proposed system's geo-mapper. This gap directly
motivates the \gls{gps}-space Kalman filter and land-mask outlier rejection stages
of the proposed pipeline, which collectively serve the same data quality function
for vision-derived coordinates that the reviewed \gls{ais} cleaning methods serve for
transponder-reported coordinates.

In \gls{ais}-to-video fusion and cross-modal track association, Chen \cite{Chen2023a}
demonstrated that bipartite matching over joint position-and-appearance cost
functions can reliably associate vision-tracked bounding boxes with \gls{ais}-reported
\gls{mmsi} identities, while Gurgel et al.\ \cite{Gurgel2010} established that \gls{ais}
records constitute a rigorous ground truth reference for validating heterogeneous
sensor-based vessel detections, a methodology directly applicable to evaluating
the \gls{gps} accuracy of the proposed system's geo-mapper. However, no reviewed work
delivers a unified real-time operator interface that combines live \gls{ais}-stream
\gls{websocket} data with a simultaneous computer-vision processing pipeline rendered
on a shared interactive map. The proposed system's \gls{ais} dashboard directly
addresses this integration gap, providing maritime operators with a single
situational awareness interface that aggregates authoritative \gls{ais} records,
\gls{mmsi}-resolved vessel identities, and computer-vision-derived track data within
one deployable application.

Together, these gaps define the original contribution of the present project:
an end-to-end maritime surveillance pipeline that integrates drone stabilisation,
\gls{yolov8} detection, \gls{deepocsort} identity-preserving tracking, \gls{gps} geo-mapping with
Kalman smoothing and outlier rejection, cold-start-aware \gls{lstm} trajectory
prediction, and a live \gls{ais} \gls{websocket} dashboard within a single deployable system
--- resolving limitations across detection, tracking, prediction, \gls{ais} data
quality, and cross-modal fusion that no single reviewed work addressed in
combination.

\section{Programming Language used}
The suggested maritime surveillance system is developed in full using Python, a high-level interpreted language which now is the de facto standard for Artificial Intelligence, Deep Learning and Computer Vision research and development. The simplicity and eloquence of Python's syntax are also driving its use in these applications — by streamlining the process of coding mathematically intensive algorithms, engineers can concentrate on designing their systems and can model the logic of the systems without having to worry about all the implementation details. However, Python is also the native interface language of almost all current, widely-used deep learning and computer vision libraries, providing seamless interoperability between all components of the proposed pipeline. It is also the most practical and productive option for an academic project this size and interdisciplinary in scope because of its extensive community support, rich documentation and its ability to be rapidly prototyped.

\section{Tools and Technologies used}
The system is based on a properly selected set of open-source libraries and frameworks that play specific roles in the pipeline. \gls{pytorch} is the backbone of deep learning, supporting dynamic building of computational graphs, automatic differentiation and native support of the \gls{cuda} libraries for the training and inference of models on \glspl{gpu}. Building on \gls{pytorch}, the Ultralytics \gls{yolov8} library offers an easy to use \gls{api} to create ship detection models, set up training and deployment parameters, and use the model for real-time inferences. Video ingestion is done with \gls{opencv}, a computer vision library that has been highly optimised for real-time frame-by-frame processing pipelines, needed for live surveillance applications, such as video ingestion. \gls{opencv} is a highly optimised computer vision library, which supports frame-by-frame processing pipelines, and therefore is indispensable for live surveillance applications, including video ingestion. 

\setlength{\parskip}{1em}

The \gls{deepocsort} tracking module is a fully integrated tracking module that fuses directly with the detection output to preserve vessel identities from frame to frame via appearance-based \gls{reid} and Kalman filter state estimation. Using \gls{numpy} and \gls{pandas}, numerical computation, array manipulation, and calculations of motion parameters (centroids) are carried out, which is fast and vectorised data handling of large trajectory datasets. Trajectory and detection visualisations are created in the video output using Matplotlib, geospatial vessel position mapping is created using Folium and interactive maps of tracked vessel positions are generated. All of the pipeline is speeded up via \gls{nvidia} \gls{cuda}, allowing the system to process the detection, tracking, and prediction tasks on real time. Testing and experimentations are done in Jupyter Notebook and Visual Studio Code, enabling iterative testing and integration of modules before rollout in the complete pipeline.


\section{Summary}

This chapter presented a comprehensive survey of the existing literature in
maritime surveillance, systematically reviewing peer-reviewed publications across
four thematic domains: deep learning-based ship detection, multi-object tracking
and \gls{reid}, vessel trajectory prediction, and \gls{ais} data integration and quality. Across all four domains,
the reviewed works demonstrated that deep learning and computer vision approaches
have produced substantial advances in detection accuracy, tracking continuity, and
trajectory forecasting capability, with the \gls{yolov8} family establishing itself as
the dominant detection framework and \gls{lstm}-based sequence models emerging as the
standard for medium-horizon trajectory prediction from \gls{ais}-derived position streams.

However, a consistent and critical gap pervades the reviewed literature: no
existing published system integrates drone-stabilised optical video detection,
appearance-based identity-preserving multi-vessel tracking, \gls{gps} geo-mapping with
coordinate noise suppression, cold-start-aware kinematic-to-\gls{lstm} trajectory
prediction, \gls{ais} data quality filtering, and a live \gls{ais} \gls{websocket} dashboard within
a single real-time deployable pipeline. Individual components have been addressed
in isolation --- outlier rejection for \gls{ais} streams \cite{Urakami2018, Wang2024a},
gap-filling via particle filter estimation \cite{Lian2019}, \gls{ais}-to-video track
association \cite{Chen2023a}, and cross-modal \gls{ais}-radar validation
\cite{Gurgel2010} --- but their integration into a unified vision-plus-\gls{ais}
surveillance system remains unaddressed in the reviewed literature.

To resolve this gap, the present project was developed using a carefully selected
software toolchain comprising Python as the primary implementation language,
\gls{pytorch} for \gls{gpu}-accelerated deep learning inference, \gls{yolov8} for per-frame vessel
localisation, \gls{deepocsort} with \gls{osnet} \gls{reid} for identity-preserving multi-object
tracking, \gls{opencv} for video processing and Lucas--Kanade optical-flow stabilisation,
\gls{numpy} for analytical motion parameter computation, and \gls{fastapi} with \gls{maplibre}
for the live \gls{ais} dashboard backend and frontend respectively. Each tool was
selected on the basis of its established performance in the reviewed literature,
its open-source availability, and its compatibility with the \gls{cuda}-accelerated
\gls{gpu} execution environment required for real-time throughput.

The research gaps identified and the technological choices justified in this
chapter feed directly into the system design and methodology presented in the
subsequent chapter, which details the architectural decisions, module interactions,
and implementation rationale behind the proposed end-to-end maritime surveillance
pipeline.